\section{Penanganan Kesalahan Leksikal}
\label{sec:ch07-error}

\subsection{Jenis Kesalahan}

Kesalahan leksikal terjadi ketika masukan tidak sesuai dengan pattern
token mana pun. Contoh umum:
\begin{itemize}[leftmargin=*]
  \item karakter ilegal, misalnya \code{@} atau \code{\#};
  \item string yang tidak ditutup;
  \item literal karakter kosong atau terlalu panjang;
  \item komentar blok yang tidak ditutup;
  \item float tanpa digit setelah titik.
\end{itemize}

Kesalahan leksikal berbeda dari kesalahan sintaksis. Lexer tidak
mengetahui struktur program, sehingga tidak dapat memberikan saran
perbaikan yang bergantung pada konteks. Namun, lexer dapat
memberikan lokasi yang tepat dan jenis masalah.

\subsection{Strategi Pemulihan}

Lexer yang baik tidak berhenti pada kesalahan pertama. Strategi yang
umum:
\begin{enumerate}[label=\arabic*.,leftmargin=*]
  \item hasilkan token \code{TOK\_ERROR};
  \item catat lokasi dan pesan;
  \item lewati satu karakter atau sampai batas aman;
  \item lanjutkan scanning.
\end{enumerate}

Untuk string yang tidak ditutup, lexer dapat melewati sampai akhir
baris. Untuk komentar blok yang tidak ditutup, lexer dapat melewati
sampai akhir berkas. Dalam kedua kasus, lexer harus tetap
menghasilkan token agar parser dapat melanjutkan analisis.

\begin{figure}[H]
\centering
\begin{tikzpicture}[node distance=6mm and 9mm]
  \node[phase, text width=2.6cm] (scan) {Scan karakter};
  \node[decision, text width=2.4cm, right=of scan] (ok) {Cocok?};
  \node[phase, text width=2.6cm, right=of ok] (token) {Hasilkan token};
  \node[phase, text width=2.6cm, below=of ok] (err) {Token error};
  \node[phase, text width=2.6cm, below=of err] (skip) {Lewati karakter};
  \draw[arrow] (scan) -- (ok);
  \draw[arrow] (ok) -- node[above, font=\small] {ya} (token);
  \draw[arrow] (ok) -- node[right, font=\small] {tidak} (err);
  \draw[arrow] (err) -- (skip);
  \draw[arrow] (skip) -| (scan);
\end{tikzpicture}
\caption{Alur pemulihan kesalahan leksikal}
\label{fig:ch07-error-flow}
\end{figure}

\subsection{Pelaporan Kesalahan}

Fungsi \code{error\_at} mencetak pesan dengan format
\code{berkas:baris:kolom: error: pesan}. Implementasinya ditunjukkan
pada \lstref{lst:ch07-error-c}.

\includecode[
  language=C,
  firstline=1,
  lastline=10,
  firstnumber=1,
  caption={Fungsi pelaporan \file{error.c}},
  label={lst:ch07-error-c}
]{../../code/subsetc/src/error.c}

Format ini mudah dibaca manusia dan mudah diurai oleh alat. Banyak
editor dan IDE dapat melompat ke lokasi kesalahan berdasarkan format
tersebut.

\subsection{Contoh Diagnostik}

\begin{table}[H]
\centering
\caption{Contoh pesan kesalahan leksikal}
\label{tab:ch07-diagnostik}
\begin{tabularx}{\textwidth}{@{}l X@{}}
\toprule
\textbf{Masukan} & \textbf{Diagnostik} \\
\midrule
\code{@} & \code{baris 1, kolom 1: karakter tidak dikenal} \\
\code{"abc} & \code{baris 1, kolom 1: string tidak ditutup} \\
\code{'ab'} & \code{baris 1, kolom 1: literal karakter terlalu panjang} \\
\code{/* abc} & \code{baris 1, kolom 1: komentar blok tidak ditutup} \\
\code{3.} & \code{baris 1, kolom 2: digit diharapkan setelah titik} \\
\bottomrule
\end{tabularx}
\end{table}

Pesan yang baik menyebut lokasi awal, jenis masalah, dan bila
mungkin tindakan yang diharapkan. Hindari pesan generik seperti
\code{error} atau \code{invalid token}.
